\chapter{Хиральное замыкание дискретной монеты}
\label{ch:v6-spectral-transition-discrete-chiral-coin-closure-gate}

\section{Минимальная заряженно-лептонная опора}

Единственный центральный блок ранга один физического пакета равен $e_R$,
но полная заряженно-лептонная опора имеет вид
\begin{equation}
 \mathcal H_{\ell}=L_L\oplus e_R\simeq\mathbb C^2\oplus\mathbb C.
\end{equation}
Проверим две возможности: обычную монету с независимым Хиггсом и более
сильную state-dependent монету, в которой хиггсовское направление строится
из самого локального состояния.

Антикруговая сверка использует результаты Тома V: пространство
калибровочно допустимых заряженно-лептонных интертвинеров одномерно, но его
амплитуда является входом конечного оператора Дирака; физический модуль
одноформ оставляет независимые относительные направления $u,d,e$.

\section{Внешний Хиггс возвращает старое юкавское ребро}

Для $H\in\mathbb C^2$ общий самосопряжённый хиральный генератор имеет вид
\begin{equation}
 K_H(y)=
 \begin{pmatrix}
  0_{2} & yH\\
  \overline y H^\dagger & 0
 \end{pmatrix}.
 \label{eq:v6-discrete-external-higgs-coin}
\end{equation}
Он ковариантен, но содержит свободную комплексную амплитуду $y$. Его
ненулевые собственные значения равны
\begin{equation}
 \pm |y|\,\|H\|,
\end{equation}
а ортогональная компонента $L_L$ остаётся нулевой модой. Нормирование
$H/\|H\|$ не определено при $H=0$. Поэтому эта ветвь является точной
дискретной записью уже известного юкавского ребра, а не новым выводом.

\section{Эндогенный составной дублет}

Пусть локальное состояние равно
\begin{equation}
 \psi=(\ell,e),\qquad \ell\in\mathbb C^2,\quad e\in\mathbb C.
\end{equation}
Минимальный билинеар с типом слабого дублета задаётся формулой
\begin{equation}
 H_{\rm eff}(\psi)=\ell\,\overline e.
 \label{eq:v6-discrete-composite-higgs}
\end{equation}
При $\ell\mapsto U\ell$ и $e\mapsto ve$ он преобразуется как
$H_{\rm eff}\mapsto U H_{\rm eff}\overline v$, то есть имеет именно
относительный тип, необходимый для хирального ребра.

Введём локальный генератор
\begin{equation}
 K(\psi)=
 \begin{pmatrix}
  0_2&H_{\rm eff}(\psi)\\
  H_{\rm eff}(\psi)^\dagger&0
 \end{pmatrix}
 \label{eq:v6-discrete-composite-generator}
\end{equation}
и обновление
\begin{equation}
 \psi^+=\exp[-i\kappa K(\psi)]\psi.
 \label{eq:v6-discrete-composite-update}
\end{equation}
Генератор самосопряжён, поэтому норма сохраняется точно. Машинный аудит
также подтверждает ковариантность генератора и полного шага относительно
случайного слабого унитарного преобразования и независимой фазы $e_R$.

Это новое содержание по сравнению с внешним $H$: направление монеты теперь
не вставляется вручную, а вычисляется из состояния. Однако коэффициент
$\kappa$ остаётся свободным.

\section{Точный закон хиральных населённостей}

Обозначим
\begin{equation}
 p=\|\ell\|^2,\qquad q=|e|^2,\qquad p+q=1.
\end{equation}
На двумерной плоскости, натянутой на $\ell$ и $e$, генератор имеет частоту
$\sqrt{pq}$. Поэтому обновление даёт точную карту
\begin{equation}
 p^+=p\cos^2(\kappa\sqrt{pq})
     +q\sin^2(\kappa\sqrt{pq}).
 \label{eq:v6-discrete-population-map}
\end{equation}
Отсюда следуют три различных инвариантных страта:
\begin{equation}
 p=0,\qquad p=1,\qquad p=\frac12.
\end{equation}
Первые два фиксированы, поскольку составной дублет исчезает на чисто
левой и чисто правой границах. Сбалансированная населенность также
сохраняется. При общих начальных данных карта может приближать $p$ к
$1/2$, но не выбирает одну из границ и не создаёт единственный хиральный
endpoint.

\begin{theorem}[Составная монета существует, но endpoint не единственен]
На $L_L\oplus e_R$ существует точная локальная, нелинейная,
нормосохраняющая и калибровочно ковариантная монета, чьё хиггсовское
направление равно билинеару $\ell\overline e$. Она устраняет необходимость
внешнего направления $H$ на уровне одного шага. Однако её динамика имеет
как минимум три инвариантных страта, не фиксирует $\kappa$ и физический
масштаб и не создаёт пространственно локализованную частицу. Поэтому R4 и
R5 остаются открытыми.
\end{theorem}

\begin{nogocriterion}
Нельзя отождествлять составной билинеар $H_{\rm eff}$ с выведенным полем
Хиггса: у него пока нет независимой кинетики, вакуумного модуля и
нормировки. Нельзя также считать $p=1/2$ массовым endpoint без проверки
пространственного профиля, спектральной щели и устойчивости.
\end{nogocriterion}

\section{Следующий гейт}

Следующий гейт
\path{version6_spectral_transition_discrete_composite_higgs_spatial_binding_gate}
должен соединить составную хиральную монету с встречным пространственным
сдвигом и проверить, существует ли конечный локализованный профиль без
внешнего потенциала.

Нулевая гипотеза: без выведенной размерной щели составная монета меняет
внутренние населённости, но пакет баллистически расплывается. Стоп-критерий:
если локализация возникает только после свободного выбора $\kappa$,
массового угла или дополнительной жёсткости, дискретная ветвь не проходит
R4--R5.

\begin{protocol}
\begin{enumerate}
 \item внешний хиггсовский интертвинер: \textbf{единственен с точностью до $y$};
 \item амплитуда $y$: \textbf{не выведена};
 \item составной дублет $\ell\overline e$: \textbf{построен};
 \item локальная калибровочная ковариантность: \textbf{точная};
 \item сохранение нормы: \textbf{точное};
 \item инвариантные страты: \textbf{$p=0,1/2,1$};
 \item единственный хиральный endpoint: \textbf{нет};
 \item пространственная локализация: \textbf{не проверена};
 \item физический масштаб и связь: \textbf{не выведены}.
\end{enumerate}
\end{protocol}

Машинный сертификат сохранён в
\path{s2t_v6_spectral_transition_discrete_chiral_coin_closure_gate_results.json}.